% I. Introducción
El dolor lumbar es la principal causa de años vividos con discapacidad en el mundo: en 2020 afectó a unos 619 millones de personas y se proyecta que alcance 843 millones de casos en 2050 \cite{gbd2023}. La resonancia magnética (RM) es la prueba de imagen de referencia para estudiar el disco intervertebral. En los cortes sagitales ponderados en T2, la pérdida de señal y de altura del disco indica degeneración, según la escala de Pfirrmann \cite{pfirrmann2001}, y la extensión del material discal más allá de su contorno define abombamientos y hernias, según la nomenclatura de consenso \cite{fardon2014}.

La lectura de estos estudios es manual, lleva tiempo y depende de la experiencia de quien la hace. Incluso con una nomenclatura estandarizada, la concordancia entre radiólogos sobre el contorno discal es solo de aceptable a moderada \cite{arana2011}. Un sistema de apoyo que clasifique de forma automática el estado de cada disco podría reducir el tiempo de lectura y hacer más homogéneos los criterios. La visión artificial ya se ha aplicado a esta tarea con buenos resultados \cite{sari2025,tan2024,sayed2025scirep,wang2025}, pero muchas cifras proceden de conjuntos de prueba pequeños, de tareas binarias o de particiones que no se declaran por paciente, y bajan al evaluar con datos externos \cite{yi2023,nigru2024}. Cuando un mismo paciente aporta imágenes a entrenamiento y a prueba, el modelo puede reconocer al paciente en lugar de la patología, y las métricas se inflan \cite{brunekreef2024,rumala2023}.

Este trabajo presenta DeepSpine, un flujo completo para clasificar discos lumbares (L1-L2 a L5-S1) en cortes sagitales T2 del \emph{Lumbar Spine MRI Dataset} \cite{sudirman2019}. El planteamiento inicial del proyecto definió cuatro clases (Normal, Degeneración, Abombamiento y Hernia) y tres preguntas de investigación:
\begin{enumerate}
  \item[P1.] \textbf{Análisis exploratorio:} ¿cómo se distribuyen los casos de cada clase en el dataset y qué diferencias visuales (brillo, forma o tamaño) permiten distinguir un disco sano de uno patológico?
  \item[P2.] \textbf{Aprendizaje clásico frente a profundo:} ¿pueden algoritmos clásicos, más sencillos y rápidos (SVM o Random Forest), alcanzar un diagnóstico confiable a partir de características de la imagen, o es necesaria una red neuronal profunda?
  \item[P3.] \textbf{Aprendizaje profundo:} ¿qué arquitectura de red convolucional logra la mayor exactitud (con una meta superior al 90\,\%) y la menor tasa de falsos negativos?
\end{enumerate}
Durante la construcción del dataset, Hernia y Abombamiento resultaron inseparables en el plano sagital y se fusionaron en una sola clase, Protrusión (Sección~\ref{sec:anotacion}); por eso las tres preguntas se responden sobre tres clases.

Las contribuciones son tres: (i) un dataset de 5\,123 ROI de 399 pacientes, construido con un proceso trazable y reproducible desde los datos públicos, que declara sus errores y limitaciones; (ii) una comparación de cuatro modelos de familias distintas (YOLOv8s-cls, ResNet-50, SVM y Random Forest) con la misma división por paciente, las mismas métricas y evaluación sobre imágenes originales; y (iii) un análisis del desbalance de clases y del sesgo por nivel vertebral que relativiza la exactitud alcanzada. La Sección~\ref{sec:relacionados} revisa los trabajos previos; la Sección~\ref{sec:metodologia} describe la construcción del dataset, la formulación del problema y los modelos; la Sección~\ref{sec:resultados} presenta y discute los resultados, y la Sección~\ref{sec:conclusiones} responde a las preguntas y expone las limitaciones y el trabajo futuro.
